% Part: computability
% Chapter: recursive-functions
% Section: partial-functions

\documentclass[../../../include/open-logic-section]{subfiles}

\begin{document}

\olfileid{cmp}{rec}{par}
\olsection{आंशिक पुनरावर्ती फलने}


पुनरावर्ती फलनांच्या व्याख्येमागील प्रेरणा समजण्यासाठी पुढील गोष्ट लक्षात
घ्या: संगणनक्षम पण आदिम पुनरावर्ती नसलेली फलने अस्तित्वात आहेत या आपल्या
सिद्धतेतून प्रत्यक्षात आणखी बरेच काही सिद्ध होते. युक्तिवाद साधा होता:
$f_0,f_1,\dots$ या फलनांचे असे प्रगणन करता येते की, $x$ आणि $y$ यांचे फलन
म्हणून $f_x(y)$ संगणनक्षम असते, एवढेच तथ्य आपण वापरले. म्हणून हा युक्तिवाद
\emph{अशा प्रकारे प्रगणन करता येणाऱ्या फलनांच्या कोणत्याही वर्गाला} लागू
होतो. यामुळे आपण अडचणीत येतो: संगणनक्षम फलनांचे स्पष्ट वर्णन आपल्याला हवे
आहे; पण संगणनक्षम फलनांच्या संग्रहाचे अशा प्रकारचे कोणतेही प्रभावी स्पष्ट
वर्णन सर्वसमावेशक असू शकत नाही!

यातून बाहेर पडण्याचा मार्ग म्हणजे \emph{अंशतः} फलनांना स्थान देणे. अंशतः
संगणनक्षम फलनांचे प्रगणन करणे \emph{शक्य आहे}, हे आपण पाहू.
\iftag{TMs}{ खरे तर, हे शक्य आहे हे आपल्याला जवळजवळ आधीच माहीत आहे,
  कारण ट्यूरिंग यंत्रांचे पद्धतशीर प्रगणन करता येते.}{} पुढे आपण आपल्या
कर्णरेषीय युक्तिवादाकडे परत येऊ आणि अंशतः फलने समाविष्ट केल्यावर तो का
चालत नाही याचा शोध घेऊ.

आता प्रश्न असा आहे: सर्व आंशिक पुनरावर्ती फलने मिळवण्यासाठी आदिम पुनरावर्ती
फलनांमध्ये काय भर घालावी लागेल? आपल्याला दोन गोष्टी कराव्या लागतील:
\begin{enumerate}
\item आदिम पुनरावर्ती फलनांची आपली व्याख्या बदलून त्यात अंशतः फलनांनाही
  स्थान द्यावे.
\item काही नवीन अंशतः फलने समाविष्ट होतील अशी काही गोष्ट व्याख्येत
  \emph{जोडावी}.
\end{enumerate}

पहिली गोष्ट सोपी आहे. पूर्वीप्रमाणे आपण शून्य, उत्तरवर्ती आणि प्रक्षेपण
यांपासून सुरुवात करू आणि संयोजन व आदिम पुनरावर्तन यांखाली बंद करू. फरक
इतकाच की, व्याख्येतील काही पदे परिभाषित नसण्याची शक्यता सामावण्यासाठी
संयोजनाची आणि आदिम पुनरावर्तनाची व्याख्या बदलावी लागते. $f$ आणि $g$ ही
अंशतः फलने असतील, तर $x$ येथे $f$ परिभाषित आहे, म्हणजे $x$ हा $f$ च्या
परिभाषाक्षेत्रात आहे, हे दर्शवण्यासाठी $f(x) \fdefined$ असे लिहू; आणि
याच्या विरुद्ध अर्थासाठी, म्हणजे $x$ येथे $f$ परिभाषित नाही हे दर्शवण्यासाठी
$f(x) \fundefined$ असे लिहू. $f(x)$ आणि $g(x)$ ही दोन्ही अपरिभाषित आहेत,
किंवा दोन्ही परिभाषित असून समान आहेत, हे दर्शवण्यासाठी $f(x) \simeq g(x)$
असे लिहू. अधिक गुंतागुंतीच्या पदांसाठीही ही चिन्हांकने वापरू. $h$ आणि
$g_0$, \dots,~$g_k$ ही सर्व अंशतः फलने असतील, तर
\[
h(g_0(\vec x),\dots,g_k(\vec x))
\]
हे पद तेव्हा आणि केवळ तेव्हाच परिभाषित असते, जेव्हा प्रत्येक $g_i$ हे
$\vec x$ येथे परिभाषित असते आणि $h$ हे $g_0(\vec x)$, \dots,~$g_k(\vec x)$
येथे परिभाषित असते, असा संकेत आपण स्वीकारू. हे समजल्यावर अंशतः फलनांसाठीच्या
संयोजनाच्या आणि आदिम पुनरावर्तनाच्या व्याख्या वरीलप्रमाणेच राहतात; फक्त
``$=$'' च्या जागी ``$\simeq$'' लिहावे लागते.

अंशतः फलने मिळवण्यासाठी आदिम पुनरावर्ती फलनांच्या व्याख्येत आपण
\emph{अपरिबद्ध शोध परिकर्मी} जोडू. $f(x,\vec z)$ हे नैसर्गिक संख्यांवरील
कोणतेही अंशतः फलन असेल, तर $\mu x \; f(x,\vec z)$ ची व्याख्या पुढीलप्रमाणे
करू:
\begin{quote}
  $f(0,\vec z), f(1,\vec z), \dots, f(x,\vec
  z)$ ही सर्व मूल्ये परिभाषित असतील आणि $f(x,\vec z) = 0$ असेल असा लघुतम
  $x$, असा $x$ अस्तित्वात असल्यास.
\end{quote}
अन्यथा $\mu x \; f(x,\vec z)$ अपरिभाषित आहे, असे समजायचे. यामुळे
$\mu x \; f(x,\vec z)$ ची एकमेव व्याख्या होते.

\begin{explain}
आपल्या व्याख्येत\iftag{TMs}{ ट्यूरिंग यंत्रांचा, किंवा}{} अल्गोरिदमचा
किंवा कोणत्याही विशिष्ट संगणन-प्रतिमानाचा संदर्भ नाही, हे लक्षात घ्या. पण
संयोजन आणि आदिम पुनरावर्तनाप्रमाणे अपरिबद्ध शोधामागेही एक क्रियात्मक,
संगणकीय सहजकल्पना आहे. अंशतः फलनाच्या संगणनक्षमतेच्या संदर्भात, ज्या
आदानांसाठी फलन अपरिभाषित असते ती आदाने संगणन न थांबणाऱ्या आदानांशी संगत
असतात. $\mu x \; f(x,\vec z)$ संगणित करण्याची प्रक्रिया अशी: शून्य मूल्य
परत येईपर्यंत $f(0,\vec z), f(1,\vec z), f(2,\vec z)$ संगणित करत राहा.
मात्र मधल्या संगणनांपैकी एखादे थांबले नाही, तर $\mu x \; f(x,\vec z)$ चे
संगणनही थांबत नाही.
\end{explain}

$R(x,\vec z)$ हा कोणताही संबंध असेल, तर $\mu x \; R(x,\vec z)$ ची व्याख्या
$\mu x \; (1 \tsub \Char{R}(x,\vec z))$ अशी करतात. दुसऱ्या शब्दांत,
$R(x,\vec z)$ सत्य असलेले $x$ चे लघुतम मूल्य $\mu x \; R(x,\vec z)$ परत
करते. म्हणून $f(x,\vec z)$ हे सर्वत्र परिभाषित फलन असेल, तर
$\mu x \; f(x,\vec z)$ हे $\mu x \; (f(x,\vec z) = 0)$ याच्याशी समान असते.
पण आपली मूळ व्याख्या अधिक सर्वसाधारण आहे, कारण ती $f(x,\vec z)$ सर्वत्र
परिभाषित नसण्याची शक्यता मान्य करते. (याउलट, संबंधाचे जातिबोधक फल नेहमीच
सर्वत्र परिभाषित असते.)

\begin{defn}
नैसर्गिक संख्यांकडून नैसर्गिक संख्यांकडे जाणाऱ्या (विविध स्थानसंख्या असलेल्या)
अंशतः फलनांचा असा लघुतम संच, ज्यात शून्य, उत्तरवर्ती आणि प्रक्षेपणे आहेत आणि
जो संयोजन, आदिम पुनरावर्तन व अपरिबद्ध शोध यांखाली बंद आहे, त्याला
\emph{आंशिक पुनरावर्ती फलनांचा} संच म्हणतात.
\end{defn}

अर्थात, काही आंशिक पुनरावर्ती फलने सर्वत्र परिभाषित असतील, म्हणजे प्रत्येक
आदानासाठी परिभाषित असतील.

\begin{defn}
\ollabel{defn:recursive-fn}
सर्वत्र परिभाषित असलेल्या आंशिक पुनरावर्ती फलनांच्या संचाला
\emph{पुनरावर्ती फलनांचा} संच म्हणतात.
\end{defn}

पुनरावर्ती फलन सर्वत्र परिभाषित आहे यावर भर देण्यासाठी त्याला कधीकधी
``सर्वत्र परिभाषित पुनरावर्ती'' असे म्हणतात.

\end{document}
